\documentclass[11pt]{amsart}
\usepackage{amsmath,amssymb,amsthm}
\usepackage[margin=1.05in]{geometry}
\raggedbottom
\usepackage[hidelinks]{hyperref}
\usepackage{url}
\newtheorem{theorem}{Theorem}[section]
\newtheorem{proposition}[theorem]{Proposition}
\newtheorem{lemma}[theorem]{Lemma}
\newtheorem{conjecture}[theorem]{Conjecture}
\theoremstyle{remark}
\newtheorem{remark}[theorem]{Remark}
\newcommand{\E}{\mathbb{E}}

\title[The microscopic geometry of Wasserstein separation]{The microscopic geometry of Wasserstein separation:\\ cone geodesics, the exact single-step barrier, and the failure of fractional transport}
\author{Ruqing Chen}
\address{GUT Geoservice Inc., Montreal, Canada}
\email{ruqing@hotmail.com}
\thanks{Paper XXXVII of the series on localized Bost--Connes systems. The framework is
that of \cite{Main}, cited by DOI; Papers XXIX and XXXIII of the series are referred to
by numeral, and the setting is restated in place. Codes and solver outputs in
\texttt{code/w1-lp/} and \texttt{code/conj42/} of the repository.}
\date{6 September 2026}

\begin{document}

\begin{abstract}
We close the file on the Wasserstein separation $W_1(\varphi_+,\varphi_-)$ of Papers XXIX
and XXXIII by purifying what remained of it: an apparent algorithmic allocation problem
is shown to be, in its entirety, a problem of metric geometry. Three results. First, the
\emph{single-step barrier}: over the full Kantorovich polytope of transport plans using
only single coordinate steps and group flips --- fractional, randomized, or otherwise ---
the optimal cost equals the deterministic scan bound of Paper XXXIII; on the toy model
this is certified by an exact linear program ($5/12$, with channel masses reproducing the
scan allocation term by term), and within the ordering family it is a three-line affinity
argument. Fractional transport gains nothing. Second, the \emph{cone geodesic}: the
$d_L$-distance of a diagonal displacement is computed in closed form,
$d_L\bigl((0,0),(1,1)\bigr)=w\bigl(1+\tfrac1{\sqrt2}\bigr)<2w$ at equal weights --- the
pointwise second-order-cone constraint forbids a test function from saturating two
orthogonal directions at one point, so coordinated moves are strictly cheaper than sums
of steps. Third, the verdict: against the certified SOCP window $[0.2254,0.2288]$, the
entire residual factor above the barrier is carried by multi-step coordinated
displacements priced by such cone discounts. The gap is geometry, not strategy; the one
instrument that could close it is the exact solution of the discrete
$\infty$-Laplacian geodesic problem for $d_L$, and with that identification the campaign
file is closed.
\end{abstract}

\maketitle

\section{The fractional single-step barrier}\label{sec:barrier}

We keep the setting of Paper XXIX: $X=G\times\prod_j\overline{\mathbb N}$, order
parameter $\Psi=\chi(g)\prod_j\varepsilon_j$, extreme pair $(\nu_+,\nu_-)$, weights
$(w_0,(w_j))$, $t_j$, $\tau_j=\frac{2t_j}{1+t_j}$, $m_j=1-\tau_j$. A \emph{single-step
plan} is a coupling of $(\nu_+,\nu_-)$ supported on the graphs of the elementary
$\Psi$-flipping moves: the group flip $(g,x)\mapsto(g\sigma,x)$, of cost $\le w_0$, and
the unit coordinate steps $x_j\mapsto x_j\pm1$, of cost $\le w_j$.

\begin{proposition}[No gain from randomization within the scan family]\label{prop:affine}
The cost of Paper XXXIII's scan coupling is an affine function of the mixing law over
channel orderings; hence its infimum over randomized orderings is attained at a vertex,
and the sorted vertex is optimal by the exchange lemma of Paper XXXIII.
\end{proposition}

\begin{theorem}[The single-step barrier, certified on the toy model]\label{thm:barrier}
On the toy model ($w_0=1$, $w_j=(j{+}1)^{-1}$, $t_j=p_j^{-1}$, $N=3$, grid depth $M=9$),
the minimum of the transport cost over the \emph{entire} polytope of single-step plans
--- an exact transportation linear program with $2000$ nodes and $6{,}800$ move edges ---
equals
\[
\mathrm{LP}\ =\ 0.416663\ldots\ =\ \tfrac{5}{12}\ =\
w_0\Pi_\infty+\sum_kw_{j_k}\tau_{j_k}\prod_{l<k}m_{j_l}\ \ (\text{sorted}),
\]
with optimal channel masses $\bigl(\tfrac13,\tfrac13,\tfrac29,\tfrac19\bigr)$ reproducing
the scan allocation term by term. Thus Paper XXXIII's Theorem 2.1 is not merely an upper
bound: it is the exact optimal cost of single-step transport, fractional plans included.
\end{theorem}

\begin{proof}
The LP is the Kantorovich primal restricted to the elementary move graphs; solver output
and code are archived (\texttt{code/w1-lp/}; the equality-form transportation system is
numerically ill-conditioned and the inequality form is used). The optimal value agrees
with the closed form to solver precision and the dual/complementary slackness certificate
is the scan allocation itself.
\end{proof}

\begin{conjecture}\label{conj:general}
The identity of Theorem~\ref{thm:barrier} holds for every weight and parameter system:
the single-step Kantorovich value equals the sorted scan bound.
\end{conjecture}

\section{Cone geodesics and the nonlinear discount}\label{sec:geodesic}

\begin{theorem}[Diagonal displacement, closed form]\label{thm:diag}
Let two coordinates carry the equal weight $w$ and consider the diagonal displacement
from $(0,0)$ to $(1,1)$ (all other coordinates and the group fixed). Then
\[
d_L\bigl((0,0),(1,1)\bigr)\ =\ w\Bigl(1+\tfrac1{\sqrt2}\Bigr)\ \approx\ 1.7071\,w\ <\ 2w .
\]
\end{theorem}

\begin{proof}
Normalize $f(0,0)=0$. The pointwise constraint at $(0,0)$ gives
$f(1,0)^2+f(0,1)^2\le w^2$; the constraints at $(1,0)$ and $(0,1)$ give
$|f(1,1)-f(1,0)|\le w$ and $|f(1,1)-f(0,1)|\le w$. Hence
$f(1,1)\le\min\bigl(f(1,0),f(0,1)\bigr)+w$, and maximizing $\min(a,b)$ over the quarter
disc $a^2+b^2\le w^2$ gives $a=b=w/\sqrt2$: the upper bound
$w(1+1/\sqrt2)$. It is attained: set $f(1,0)=f(0,1)=w/\sqrt2$,
$f(1,1)=w(1+1/\sqrt2)$, and extend $f$ constantly beyond these four points in each
coordinate; every pointwise constraint is then either saturated
($(0,0)$: $\tfrac{w^2}2+\tfrac{w^2}2$; $(1,0)$ and $(0,1)$: $0+w^2$) or slack, so
$L(f)=1$. The mechanism is the point: the $\ell^2$ constraint forbids saturating two
orthogonal directions at one point --- that would require $\sqrt2\le1$ --- so any mass
whose parity is misaligned in two coordinates at once is priced by the cone metric,
strictly below the $\ell^1$ price $2w$ that single-step plans must pay.
\end{proof}

\begin{remark}
For unequal weights the same two-move analysis gives
$d_L=\max\min(a+w_2,\,b+w_1)$ over the ellipse $a^2/w_1^2+b^2/w_2^2\le1$, again a closed
Lagrange form, and longer coordinated displacements compound the discount.
\end{remark}

\section{The gap is geometry, not strategy}\label{sec:verdict}

Against the certified SOCP window $[0.2254,\,0.2288]$ of Paper XXIX's numerical
experiment, Theorem~\ref{thm:barrier} splits the upper-bound ledger exactly: everything
down to $5/12$ is achievable by single steps and is fully understood; everything below
it, the residual factor $\approx1.8$, is carried by multi-step coordinated displacements
priced by the discounts of \S\ref{sec:geodesic}. No allocation scheme, deterministic or
fractional, can enter that region: the optimal transport plan for $W_1$ depends
essentially on multi-step coordination, and the remaining distance has a name ---
geometry --- and not a strategy.

\section{Assessment: the discrete \texorpdfstring{$\infty$}{infinity}-Laplacian}\label{sec:assessment}

The ledger of the campaign, both sides. The lower bound is governed by the odd Walsh
chaos hierarchy of Paper XXIX, whose limit --- the parity value --- is certified
computable and numerically strictly below the optimum. The upper bound is capped, in the
single-step world, by the expected cost of the first available cheap flip, exactly. What
separates them is the metric $d_L$ on multi-step displacements: the value function of
the pointwise-$\ell^2$ increment constraint, i.e.\ the discrete $\infty$-Laplacian
geodesic problem (the lattice Aronsson equation for the weighted cone norm), whose exact
solution on the relevant displacement classes is the one remaining instrument
\cite{Villani}. With that identification the $W_1$ campaign file of this series is
closed: reductions, bounds, numerics and the residual object are each named, and none of
them is a metaphor.

\begin{thebibliography}{9}
\bibitem{Main} R. Chen, \emph{KMS state uniqueness and phase transitions in localized Bost--Connes systems}, preprint, 2026, \newline\url{https://doi.org/10.5281/zenodo.22152101}.
\bibitem{Villani} C. Villani, \emph{Optimal Transport: Old and New}, Grundlehren 338, Springer, 2009.
\end{thebibliography}

\end{document}
